% Pursuing Stacks (cotes 134-2 to 134-8): archive cote 134-5, site page 87 (PDF page 88), typescript page 274.
% typescript. Draft compared: ps4.tex:4458-4524.
% AI-transcribed by Claude Opus 5.5 (claude-opus-5-5, adaptive thinking, effort high), Message Batches API, batch ps2 (msgbatch_01NVwr1Yuo8qJf5rykeGBXPK), 2026-11-12; not yet checked by eye.
\scanpage{134-5/87}
\mspage{274}

29.6.

84. Two days have passed without writing \over{the}{any} notes. Much of the
time I spent on writing mathematical letters - one pretty long one to
Gerd Falting\hand{s}, who (on my request) had sent me preprints of his \xout{late}
recent work,\inserted{notably} on the Tate conjectures for abelian varieties and on the
Mordell conjecture, and had expressed interest hearing about some ideas
and conjectures on ``anabelian algebraic geometry''. I had been impressed,
from a glance upon the last\inserted{of his} manuscript\unsure{s}, to see three\inserted{key} conjectures \xout{which}
proved in about forty pages, while they were being considered as quite
out of reach by the people supposed to know. Some ``anabelian'' conjectures
of mine are closely related to the Tate and Mordell conjectures just pro-%
ved by Falting - Deligne had pointed out to me about two years ago that
a certain fixed-point conjecture (which I like to view as the basic conjec-%
ture at present in the anabelian program) implied Mordell's, so why loose
one's time on it ! I have the feeling Falting is the kind of chap who
may become interested in things which are supposed to be too far off to
be worth looking at, that's why I had written him a few words, under
the moment's inspiration. - Another letter, not quite as long, was an
answer to a very long and patient letter of Tim Porter, te\over{a}{l}ling me about
a number of things which have been done by homotopy people,\inserted{and} which I was
\inserted{of course} wholly ignorant of. His letter has been the first echo I got from some-%
one who read part of the notes on ``Pursuing stacks'' , and I was glad he
could make some sense of what he read so far, and conversely - that not
all he was telling me was going wholly ``above my head'' !

Apart from this and the (not unpleasant !) daily routine, I spent
a fair bunch of hours on scratchwork, centering around trying to figure
out the right notion of morphism for asphericity structures. It was a sur-%
prise that the notion should be so reticent for revealing itself - as a
matter of fact, I am not quite sure yet if I got the right notion, in
sufficient generality I mean. Time will tell - for the time being,
the notion of morphism I have to offer, while maybe too restrictive,
looks really seducing, because it parallels so perfectly the formalism
of aspheric functors $i: A \hand{\longrightarrow} M$ and the corresponding (would-be modeli-%
zing) functor $i^{\hand{*}}: M \hand{\to} A^{\wedge}$. Again, it has been a surprise, right now,
that after uncounted hours of unconvincing efforts today and yesterday,
and almost wholly unrelated to these, this pretty set up would come out
within ten minutes reflection !

\margin{87}
